\input{../tex-inputs/latex-leseansicht-vorspann}

\section[Arthur Schnitzler an Gustav Schwarzkopf, 21. 9. 1916]{L04409 Arthur Schnitzler an Gustav Schwarzkopf, 21. 9. 1916}
\nopagebreak\mylabel{L04409v}
\rehead{ }\normalsize\beginnumbering\briefempfaengerindex{Schwarzkopf, Gustav@\textsc{Schwarzkopf, Gustav}!zzzSchnitzler, Olga@\emph{von Olga Schnitzler}!1916-09-211@{21. 9. 1916}|(be}\briefempfaengerindex{Schwarzkopf, Gustav@\textsc{Schwarzkopf, Gustav}!zzzSchnitzler, Arthur@\emph{von Arthur Schnitzler}!1916-09-211@{21. 9. 1916}|(be}
\toendnotes[C]{\smallbreak\pagebreak[2]}
\correspDesc{Versand  durch Arthur Schnitzler, Olga Schnitzler am 21. 9. 1916 in Altaussee
\newline{}Übermittlung  am 21. 9. 1916 in Altaussee
\newline{}Erhalt  durch Gustav Schwarzkopf im Zeitraum [22. 9. 1916 – 26. 9. 1916?] in Wien}\toendnotes[C]{\smallbreak}
\Standort{DLA, A:Schnitzler, HS.1985.1.1897.}
\physDesc{Bildpostkarte, 692 Zeichen
\newline{}Handschrift: Bleistift, deutsche Kurrent
\newline{}Versand: Stempel: »\nobreak{}\oindex{Altaussee@\textbf{Altaussee}|pwk}Alt Aussee, 21. IX. 16\nobreak{}«.  }\toendnotes[C]{\smallbreak}\pstart{}{\pb}Herrn Guſtav Schwarzkopf\pend{}\pstart{}Wien I\oindex{I., Innere Stadt@\textbf{I., Innere Stadt}|pw}\pend{}\pstart{}Tiefer Graben 17\oindex{Wien@\textbf{Wien}!I., Innere Stadt@\textbf{I., Innere Stadt}!Tiefer Graben 17@\textbf{Tiefer Graben 17}|pw}\pend{}{\bigskip}
\pstart
           \centering{}{\pb}\textcolor{gray}{\textbf{Salzkammergut\oindex{Salzkammergut@\textbf{Salzkammergut}|pw}. Bad Aussee\oindex{Bad Aussee@\textbf{Bad Aussee}|pw} und Umgebung aus der Vogelschau.}}\pend
           
\pstart
           \centering{}\textcolor{gray}{\textbf{Loser\oindex{Loser@\textbf{Loser}|pw}\hspace*{1.5em}Trisselwand\oindex{Trisselwand@\textbf{Trisselwand}|pw}{ }Backstein\oindex{Backenstein@\textbf{Backenstein}|pw}\hspace*{1.5em}Gössl\oindex{Gössl@\textbf{Gössl}|pw}\hspace*{1.5em}z. Toplitz-See\oindex{Toplitzsee@\textbf{Toplitzsee}|pw}}}\pend
           
\pstart
           \centering{}\textcolor{gray}{\textbf{Seewiese\oindex{Seewiese [Altaussee]@\textbf{Seewiese [Altaussee]}|pw}\hspace*{1.5em}Klamm Kgl.\oindex{Klammkogel@\textbf{Klammkogel}|pw}\hspace*{1.5em}Schraml\oindex{Schramlgut@\textbf{Schramlgut}|pw}\hspace*{1.5em}Grundl-See\oindex{Grundlsee [See]@\textbf{Grundlsee [See]}|pw}}}\pend
           
\pstart
           \centering{}\textcolor{gray}{\textbf{Fischerndorf\oindex{Fischerndorf@\textbf{Fischerndorf}|pw}\hspace*{1.5em}Altausseer-See\oindex{Altausseer See@\textbf{Altausseer See}|pw}\hspace*{1.5em}Tressenstein\oindex{Tressenstein@\textbf{Tressenstein}|pw}\hspace*{1.5em}Klause\oindex{Seeklause@\textbf{Seeklause}|pw}}}\pend
           
\pstart
           \centering{}\textcolor{gray}{\textbf{Altaussee\oindex{Altaussee@\textbf{Altaussee}|pw}\hspace*{1.5em}Praunfalk\oindex{Praunfalk@\textbf{Praunfalk}|pw}\hspace*{1.5em}H. Elisabeth\oindex{Bade-Hôtel Elisabeth@\textbf{Bade-Hôtel Elisabeth}|pw}\hspace*{3em}St. Leonhard\oindex{Leonhardkirche Bad Aussee@\textbf{Leonhardkirche Bad Aussee}|pw}}}\pend
           
\pstart
           \centering{}\textcolor{gray}{\textbf{Wasner\oindex{Wasnerin@\textbf{Wasnerin}|pw}\hspace*{1.5em}Bad Aussee\oindex{Bad Aussee@\textbf{Bad Aussee}|pw}\hspace*{1.5em}Hohe Radling\oindex{Hohe Radling@\textbf{Hohe Radling}|pw}}}\pend
           
\pstart
           \centering{}\textcolor{gray}{\textbf{Sarstein\oindex{Sarstein@\textbf{Sarstein}|pw}\hspace*{1.5em}\hspace*{1.5em}Unt.
                     Kainisch\oindex{Unterkainisch@\textbf{Unterkainisch}|pw}{ } n. Stainach\oindex{Stainach@\textbf{Stainach}|pw}}}\pend
           
\pstart
           \centering{}\textcolor{gray}{\textbf{Bhf. Aussee\oindex{Bahnhof Bad Aussee@\textbf{Bahnhof Bad Aussee}|pw}\hspace*{3em}Koppenthal\oindex{Koppental@\textbf{Koppental}|pw}}}\pend
           
\pstart
           \centering{}\textcolor{gray}{\textbf{n. Ischl\oindex{Bad Ischl@\textbf{Bad Ischl}|pw}}}\pend
           \vspace{1em}
\pstart
           \raggedleft{}{\pb}21. 9. 16\pend
           \vspace{0.5em}
\pstart
           Noch immer, lieber Guſtav, ſind wir hier\oindex{Altaussee@\textbf{Altaussee}|pw}, bei miſerabelſtem Wetter, auch liegt der Schnee { }ſchon bis tief
               herunter, aber die Verwickeltheit unſrer häuslichen Zuſtände (\label{K_L04409-1v}\edtext{neues Fräulein\pwindex{Pfister, Angela-Julianna 22.\,2.\,1885 Aletshausen – 19.\,6.\,1917 Probstdeuben@\textsc{Pfister, Angela-Julianna} (22.\,2.\,1885 Aletshausen – 19.\,6.\,1917 Probstdeuben)|pwv}}{\lemma{\textnormal{\emph{neues Fräulein}}}\Cendnote{\textnormal{Vgl. A. S.: \emph{Tagebuch}, 15. 9. 1916.}}}\label{K_L04409-1}
               u. ä.) veranlaſst uns die Reiſe nach Hauſe bis gegen Ende d. M. aufzuſchieben{[}.{]} Übrigens läſst es ſich in
               der Villa\oindex{Seevilla [Altaussee]@\textbf{Seevilla [Altaussee]}|pwv}, die gut heizbar
               iſt, ganz behaglich leben (was freilich nicht der Sinn des So{\geminationm}eraufenthalts iſt – das Heizen mein ich, – nicht d 
               Behaglichkeit. Heini\pwindex{Schnitzler, Heinrich 9.\,8.\,1902 Hinterbrühl – 12.\,7.\,1982 Wien@\textsc{Schnitzler, Heinrich} (9.\,8.\,1902 Hinterbrühl – 12.\,7.\,1982 Wien)|pw} iſt ſchon ſeit 8 Tagen
               (mit Fingi\pwindex{Jacobus, Elisabeth 4.\,8.\,1877 Trier – 10.\,1.\,1963 Wien@\textsc{Jacobus, Elisabeth} (4.\,8.\,1877 Trier – 10.\,1.\,1963 Wien)|pw}) in Wien\oindex{Wien@\textbf{Wien}|pw}; hat er Ihnen ſchon {\pb}telephonirt? Muſik
               geht auch nicht immer, ich{ }ſchreibe fleißig, 4 händig. Aber die Welt iſt gar nicht{ }ſchön\pend
           
\pstart
           herzlichſt Ihr{\\[\baselineskip]}\spacefill\mbox{A.}\pend
           \leftskip=0em{}
\pstart
           \noindent{}Olga grüßt! –\pend
           \selectlanguage{ngerman}\endnumbering\briefempfaengerindex{Schwarzkopf, Gustav@\textsc{Schwarzkopf, Gustav}!zzzSchnitzler, Olga@\emph{von Olga Schnitzler}!1916-09-211@{21. 9. 1916}|)be}\briefempfaengerindex{Schwarzkopf, Gustav@\textsc{Schwarzkopf, Gustav}!zzzSchnitzler, Arthur@\emph{von Arthur Schnitzler}!1916-09-211@{21. 9. 1916}|)be}\mylabel{L04409h}
\begin{anhang}
\end{anhang}\newcommand{\dateiname}{L04409}\newcommand{\titel}{Arthur Schnitzler an Gustav Schwarzkopf, 21. 9. 1916}\newcommand{\editorInnen}{Herausgegeben von Jahnke, SelmaMüller, Martin Anton}\input{../tex-inputs/latex-leseansicht-abspann}
